% =====================================================================
% ANNEXES — RÉSERVE (contenu retiré du rapport pour respecter la
% contrainte de 100 pages totales, annexes incluses)
% Fichier NON injecté dans le rapport compilé. Conservé pour
% réutilisation future (rapport étendu, soutenance, documentation
% interne).
% =====================================================================

% =====================================================================
% Ex-Annexe B~--- Structure complète du bundle Databricks
% =====================================================================

Cette annexe détaille la structure complète d'un bundle Databricks généré par ADMP~AI~Studio, à titre d'exemple pour un projet de migration standard composé de trois unités de migration (\texttt{CUS01}, \texttt{ACC00}, \texttt{APD}).

\subsection*{Arborescence du bundle}

\begin{tcolorbox}[colback=accenturegray!40, colframe=accenturedark, boxrule=0.5pt, arc=3pt,
  left=8pt, right=8pt, top=5pt, bottom=5pt, fontupper=\small\ttfamily]
poc\_\textit{project}\_YYYYMMDD\_bundle/\newline
\quad poc\_\textit{project}\_YYYYMMDD/\newline
\quad\quad runner.py\newline
\quad\quad validate.py\newline
\quad\quad MANIFEST.json\newline
\quad\quad units/\newline
\quad\quad\quad mig\_cus01.py\newline
\quad\quad\quad mig\_acc00.py\newline
\quad\quad\quad mig\_apd.py\newline
\quad\quad framework/\newline
\quad\quad\quad unit\_base.py\newline
\quad\quad config/\newline
\quad\quad\quad parameters.py\newline
\quad\quad\quad transco.py\newline
\quad\quad source\_ddl/\newline
\quad\quad\quad create\_tables.py\newline
\quad\quad\quad generate.py\newline
\quad\quad reports/\newline
\quad\quad\quad generate\_reports.py\newline
\quad\quad\quad integrity\_checks.json\newline
poc\_\textit{project}\_YYYYMMDD.whl
\end{tcolorbox}

\subsection*{Rôle de chaque composant}

\begin{table}[H]
\centering
\caption{Description des composants du bundle Databricks}
\label{tab:bundle-detail}
\begin{tabularx}{\textwidth}{llX}
\toprule
\textbf{Fichier} & \textbf{Type} & \textbf{Rôle} \\
\midrule
\texttt{runner.py} & Orchestrateur &
  Point d'entrée principal~; exécute les unités dans l'ordre configuré~; passe les paramètres run-time (catalog, schéma, run\_id) \\
\addlinespace
\texttt{validate.py} & Validation pré-run &
  Vérifie l'existence des tables cibles dans Unity Catalog~; contrôle la cohérence des paramètres avant toute écriture \\
\addlinespace
\texttt{MANIFEST.json} & Méta-données &
  Contient la version du bundle, les unités incluses, la date de génération et la version du compilateur \\
\addlinespace
\texttt{units/mig\_*.py} & Unités PySpark &
  Chaque fichier contient la logique de migration d'une entité~: DataFrame Spark, UNION ALL des branches, LEFT JOIN transco, INSERT INTO Delta \\
\addlinespace
\texttt{framework/unit\_base.py} & Classe de base &
  Fournit les mécanismes de logging (compteurs lignes lues/écrites/rejetées), gestion des erreurs et accès aux paramètres \\
\addlinespace
\texttt{config/parameters.py} & Paramètres &
  Dictionnaire typé des paramètres métier et des références catalog/schéma, modifiable avant chaque run \\
\addlinespace
\texttt{config/transco.py} & Transcodification &
  Charge les tables de transco depuis les fichiers de référence et les diffuse en broadcast Spark pour optimiser les jointures \\
\addlinespace
\texttt{source\_ddl/create\_tables.py} & DDL source &
  Recrée les tables sources en mémoire (tables Delta temporaires) pour permettre une exécution offline sans connexion aux systèmes sources \\
\addlinespace
\texttt{source\_ddl/generate.py} & Données de test &
  Génère des données de test synthétiques pour les tables sources~; utilisé dans les environnements de validation (UC-20) \\
\addlinespace
\texttt{reports/generate\_reports.py} & Rapports Excel &
  Génère les rapports de contrôle d'intégrité référentielle par unité et le rapport consolidé global \\
\addlinespace
\texttt{reports/integrity\_checks.json} & Définitions &
  Contient la définition des contrôles d'intégrité à exécuter~: clés primaires, clés étrangères, unicité, règles de gestion \\
\bottomrule
\end{tabularx}
\end{table}

% =====================================================================
% Ex-Annexe C~--- Modèle de données PostgreSQL
% =====================================================================

Cette annexe présente les principales tables de la base de données PostgreSQL d'ADMP~AI~Studio, telles que définies dans les migrations Flyway du dépôt. Les colonnes de verrouillage de section (checkout) sont détaillées car elles constituent un mécanisme architectural clé.

\begin{table}[H]
\centering
\caption{Table \texttt{migration\_units}~--- colonnes de verrouillage de section}
\label{tab:migration-units-lock}
\begin{tabularx}{\textwidth}{llX}
\toprule
\textbf{Colonne} & \textbf{Type} & \textbf{Description} \\
\midrule
\texttt{id} & UUID & Identifiant unique de l'unité de migration \\
\texttt{project\_id} & UUID & Référence au projet parent (FK) \\
\texttt{name} & VARCHAR & Nom de l'unité (ex.~\texttt{CUS01}, \texttt{ACC00}) \\
\texttt{data\_structure\_locked\_by} & UUID & ID utilisateur détenant le verrou sur la section Data Structure~; \texttt{NULL} si section libre \\
\texttt{record\_creation\_locked\_by} & UUID & ID utilisateur détenant le verrou sur la section RC \\
\texttt{data\_population\_locked\_by} & UUID & ID utilisateur détenant le verrou sur la section DP \\
\texttt{checkout\_baseline} & INTEGER & Numéro de version de référence au moment du checkout~; utilisé pour détecter les conflits \\
\texttt{created\_at} & TIMESTAMPTZ & Date de création de l'unité \\
\texttt{updated\_at} & TIMESTAMPTZ & Date de dernière modification \\
\bottomrule
\end{tabularx}
\end{table}

Le mécanisme de verrouillage garantit l'exclusivité d'édition par section~: chaque colonne \texttt{*\_locked\_by} peut être \texttt{NULL} (section libre) ou contenir l'UUID de l'utilisateur qui détient le verrou. La colonne \texttt{checkout\_baseline} permet de détecter les modifications concurrentes~: si la version courante de l'unité diffère de la baseline au moment du check-in, un conflit est signalé et une résolution manuelle est demandée.
